% 자동 생성: verification/agent_colloc_20261015.json
\begin{widetable}
\centering
\caption{주요 연어의 비중 변화와 연구 대상의 구체화. 비율의 분모는 각 연도의 전체 집합이며, 2026년은 1$\sim$9월이다.}
\label{tab:colloc-direction}
\setlength{\tabcolsep}{5pt}
\begin{tabular}{@{}>{\raggedright\arraybackslash}p{0.20\textwidth}>{\raggedright\arraybackslash}p{0.20\textwidth}>{\centering\arraybackslash}p{0.18\textwidth}>{\raggedright\arraybackslash}p{0.36\textwidth}@{}}
\toprule
구체화되는 연구 대상 & 연어 & 2023년 $\rightarrow$ 2026년 & 읽을 수 있는 변화 \\
\midrule
수행하는 작업 & coding agent & 0.0\% $\rightarrow$ 4.2\% & 코딩처럼 수행 목적을 명시한 에이전트 연구의 확대 \\
\addlinespace[3pt]
작업을 연결하는 절차 & agent workflow & 0.0\% $\rightarrow$ 1.3\% & 여러 단계의 작업을 어떻게 구성하고 연결할지에 관한 관심 \\
\addlinespace[3pt]
작업 중 정보를 유지하는 방식 & agent memory & 0.2\% $\rightarrow$ 1.6\% & 기억을 에이전트의 구성 요소로 다루는 연구의 확대 \\
\addlinespace[3pt]
수행 결과를 비교하는 기준 & agent benchmark & 0.3\% $\rightarrow$ 1.4\% & 에이전트의 수행 능력을 평가하는 연구의 확대 \\
\bottomrule
\end{tabular}
\end{widetable}
